% !TeX root = ../初等数学笔记.tex
% 由主文档通过 \input 载入；请编译主文档。
\setcounter{chapter}{1} % 保留原书章号。
\chapter{分层自测与专题阅读}
\label{app:em-study}
学习检查应同时覆盖概念、计算、证明和解释。能够识别一个熟悉公式，并不等于能够说明它的条件；能够跟随一个证明，也不等于能够独立补足关键步骤。

\section{分阶段自测与复习路径}
第一阶段检查集合包含与元素关系、充分必要条件、绝对值、整除及定义域。第二阶段检查方程变形是否等价、函数是否单射、参数何时退化、数列公式是否覆盖初值。第三阶段检查不等式等号能否达到、计数是否重复、有限证明中所选对象是否存在。

几何阶段检查图形中的关系是否来自条件、全等相似的对应是否一致、轨迹是否双向验证、作图是否遗漏退化情况；三角与坐标阶段检查角度单位、分支、法向量和曲线参数约定。概率统计阶段检查模型是否等可能、独立性是否给出、样本结论能否扩展到总体。遇到概念性错误应回到定义，遇到运算性错误应逐步验算，遇到证明性错误应寻找缺失依据或反例。

\section{综合自测及答案要点}
\begin{enumerate}
\item 解 $\dfrac{x^2-1}{x-1}\ge0$。答：$[-1,1)\cup(1,+\infty)$，先保留 $x\ne1$ 再化简。
\item 正数 $x,y$ 满足 $x+y=10$，求 $x^2+y^2$ 的最小值。答：$50$，由 $(x-y)^2\ge0$，在 $x=y=5$ 取得。
\item 求从 $7$ 个不同对象中选出 $3$ 个的方法数，并求同时选中两个指定对象的概率，假设所有三元子集等可能。答：总数 $35$，有利数 $5$，概率 $1/7$。
\item 三角形三边为 $5,5,6$，求面积、内切圆和外接圆半径。答：高为 $4$，面积 $12$；由 $S=rs$ 得内切半径 $3/2$，由 $S=abc/(4R)$ 得外接半径 $25/8$。前一关系来自内心分成的三个三角形，后一关系结合正弦定理与面积式。
\item 解复数方程 $z^2=-4$。答：$\pm2i$；实数范围内则无解，先确定数系。
\item 给出一个样本均值正确但总体解释可能错误的情形。要点：例如仅调查自愿参加者却把结果扩展到所有对象，错误在抽样代表性而非均值公式。
\end{enumerate}

\section{现有专题与进一步阅读}
现有均值、初等不等式与 Jensen 不等式专题适合在第\ref{chap:em-inequality-methods}章之后阅读，其中涉及积分或一般随机变量的推广需补足相应分析与概率基础。裂项公式和双阶乘专题分别衔接第\ref{chap:em-sequences}、\ref{chap:em-counting}章；三角公式专题衔接第\ref{chap:em-triangle-solving}章；圆与球的度量专题可用于进一步追踪长度、面积和体积逼近的严格依据。

继续学习时，数学分析承接实数完备性、数列与函数，线性代数承接方程组和向量，初等数论与抽象代数深化整除及代数结构，概率论与数理统计扩展有限模型和抽样推断。已有知识的价值在于明确哪些工具已经证明、哪些条件不可省略，以及新问题究竟超出了现有理论的哪一步。

